\documentclass[12pt]{article}
\usepackage[utf8]{inputenc}
\usepackage[T1]{fontenc}
\usepackage[margin=1in]{geometry}
\usepackage{setspace}
\usepackage{xcolor}
\usepackage[colorlinks=true, urlcolor=blue!60!black, linkcolor=blue!60!black]{hyperref}
\setlength{\parindent}{2em}
\setlength{\parskip}{1.25em}
\renewcommand{\baselinestretch}{1.0}
\begin{document}
\begin{center}
{\LARGE \textit{Green v. County School Board of New Kent County}, 391 U.S. 430 (1968)\par}
\vspace{0.8em}
{\large Case Brief\par}
\end{center}
\vspace{1em}

\section{Memory Jogger}

School boards that operated state-mandated dual school systems cannot satisfy their affirmative Fourteenth Amendment duty to eliminate segregation ``root and branch'' by adopting a ``freedom-of-choice'' plan that leaves racial identifiability intact when more effective alternatives such as geographic zoning are readily available.

\section{Detailed Case Facts}

New Kent County, a rural county in eastern Virginia, had a population of approximately 4,500 residents in the 1960s, about half of whom were Black (\href{https://supreme.justia.com/cases/federal/us/391/430/}{391 U.S. at 432}; \href{https://www.courtlistener.com/opinion/107705/green-v-county-school-board-of-new-kent-county/}{CourtListener}). Unlike many urban districts, New Kent County had no residential segregation: Black and white residents lived scattered throughout the entire county. The county school system served roughly 1,300 pupils---approximately 740 Black students and 550 white students---and operated only two combined elementary-and-high schools: the New Kent School on the east side of the county for white students and the George W. Watkins School on the west side for Black students (\href{https://supreme.justia.com/cases/federal/us/391/430/}{391 U.S. at 432}). There were no geographic attendance zones; instead, each school served the entire county, and 21 school buses---11 serving Watkins and 10 serving New Kent---traveled overlapping routes across the county to transport students to their racially designated schools (\href{https://supreme.justia.com/cases/federal/us/391/430/}{391 U.S. at 432}).

This dual system had been established under Virginia constitutional and statutory mandates (Va. Const., Art. IX, \S~140 (1902); Va. Code \S~22-221 (1950)) that the Supreme Court had invalidated in \textit{Davis v. County School Board of Prince Edward County}, decided alongside \textit{Brown v. Board of Education}, 347 U.S. 483 (1954) (\textit{Brown I}) (\href{https://supreme.justia.com/cases/federal/us/391/430/}{391 U.S. at 432--33}). Following \textit{Brown I} and \textit{Brown II}, 349 U.S. 294 (1955), the New Kent County School Board continued operating completely segregated schools under Virginia's ``Massive Resistance'' statutes, including the Pupil Placement Act (Va. Code \S~22-232.1 \textit{et seq.} (1964), not repealed until 1966) (\href{https://encyclopediavirginia.org/entries/green-charles-c-et-al-v-county-school-board-of-new-kent-county-virginia/}{Encyclopedia Virginia}; \href{https://supreme.justia.com/cases/federal/us/391/430/}{391 U.S. at 433}). Under that Act, a State Pupil Placement Board automatically reassigned children to their previous schools unless they applied for and received state approval to transfer. By September 1964---a decade after \textit{Brown I}---not a single Black student had applied to attend New Kent School, and no white student had applied to attend Watkins School (\href{https://supreme.justia.com/cases/federal/us/391/430/}{391 U.S. at 433}).

In March 1965, Charles C. Green and other Black parents and students sued the School Board for injunctive relief. Five months later, on August 2, 1965, in order to remain eligible for federal financial assistance under Title VI of the Civil Rights Act of 1964 (42 U.S.C. \S~2000d) and Department of Health, Education, and Welfare guidelines, the Board adopted a ``freedom-of-choice'' plan (\href{https://supreme.justia.com/cases/federal/us/391/430/}{391 U.S. at 433 \& n.2}). The plan allowed students---except those entering the first and eighth grades, who had to make an affirmative choice---to choose annually between the two schools; students who did not submit a choice were automatically reassigned to the school they had previously attended (\href{https://supreme.justia.com/cases/federal/us/391/430/}{391 U.S. at 433--34}). Over three years of operation (1965--1967), not a single white student chose to attend Watkins School, and although 35 Black students enrolled at New Kent in 1965, 111 in 1966, and 115 in 1967, 85\% of the county's Black students still attended the all-Black Watkins School (\href{https://supreme.justia.com/cases/federal/us/391/430/}{391 U.S. at 441}; \href{https://www.oyez.org/cases/1967/695}{Oyez}).

\section{Procedural History}

\begin{itemize}
\item \textbf{United States District Court for the Eastern District of Virginia (1965--1966):} Petitioners filed suit in March 1965 seeking an injunction against continued maintenance of a racially segregated school system. After the Board adopted its freedom-of-choice plan on August 2, 1965, the District Court denied the injunction and granted the Board leave to submit an amendment addressing nondiscriminatory employment and assignment of faculty and staff. On June 2, 1966, following a May 17, 1966 memorandum, the District Court approved the amended freedom-of-choice plan (\href{https://supreme.justia.com/cases/federal/us/391/430/}{391 U.S. at 432, 434}).
\item \textbf{United States Court of Appeals for the Fourth Circuit (1967):} Sitting en banc, the Fourth Circuit decided the case per curiam (\textit{Green v. County School Board of New Kent County}, 382 F.2d 338 (4th Cir. 1967)) on the basis of its companion decision in \textit{Bowman v. County School Board of Charles City County}, 382 F.2d 326 (4th Cir. 1967) (\href{https://supreme.justia.com/cases/federal/us/391/430/}{391 U.S. at 434 \& n.3}). The Fourth Circuit affirmed approval of the freedom-of-choice provisions but remanded for a more specific and comprehensive order governing faculty desegregation, including an objective timetable. Judges Sobeloff and Winter concurred in the faculty remand but disagreed with upholding freedom of choice, urging periodic judicial evaluation of its effectiveness (\href{https://supreme.justia.com/cases/federal/us/391/430/}{391 U.S. at 434--35}).
\item \textbf{United States Supreme Court (1967--1968):} The Court granted certiorari on December 11, 1967 (389 U.S. 1003), heard oral argument on April 3, 1968 (alongside companion cases \href{https://supreme.justia.com/cases/federal/us/391/443/}{\textit{Raney v. Board of Education}}, 391 U.S. 443 (1968), and \href{https://caselaw.findlaw.com/court/us-supreme-court/391/450.html}{\textit{Monroe v. Board of Commissioners}}, 391 U.S. 450 (1968)), and issued its decision on May 27, 1968, vacating the Court of Appeals' judgment insofar as it affirmed the District Court and remanding for further proceedings (\href{https://www.oyez.org/cases/1967/695}{Oyez}; \href{https://supreme.justia.com/cases/federal/us/391/430/}{391 U.S. at 442}).
\end{itemize}

\section{Judicial Votes}

\textbf{Vote:} Unanimous, 9--0 (\href{https://www.oyez.org/cases/1967/695}{Oyez}; \href{https://supreme.justia.com/cases/federal/us/391/430/}{Justia}).

\begin{itemize}
\item \textbf{Majority Opinion:} Justice William J. Brennan, Jr., joined by Chief Justice Earl Warren and Justices Hugo L. Black, William O. Douglas, John M. Harlan II, Potter Stewart, Byron R. White, Abe Fortas, and Thurgood Marshall.
\item \textbf{Concurrences:} None.
\item \textbf{Dissents:} None.
\end{itemize}

\section{Holding}

\textbf{Judgment:} The judgment of the Court of Appeals for the Fourth Circuit was vacated insofar as it affirmed the District Court's approval of the freedom-of-choice plan, and the case was remanded to the District Court for further proceedings requiring the School Board to formulate a new plan (\href{https://supreme.justia.com/cases/federal/us/391/430/}{391 U.S. at 442}).

\textbf{Rule Announced:} A school board that operated a state-compelled dual school system at the time of \textit{Brown I} bears an affirmative Fourteenth Amendment duty to take whatever steps are necessary to convert promptly to a unitary, nonracial system in which racial discrimination is eliminated ``root and branch'' across every facet of school operations---student assignment, faculty, staff, transportation, extracurricular activities, and facilities (later known as the \textit{Green} factors) (\href{https://supreme.justia.com/cases/federal/us/391/430/}{391 U.S. at 435, 437--38}). The board's burden is to produce a plan that ``promises realistically to work, and promises realistically to work now'' (\href{https://supreme.justia.com/cases/federal/us/391/430/}{391 U.S. at 439}). While a freedom-of-choice plan is not unconstitutional per se, it is unacceptable whenever it fails to dismantle the dual system and reasonably available alternatives---such as geographic zoning or grade consolidation---promise speedier and more effective conversion to a unitary system (\href{https://supreme.justia.com/cases/federal/us/391/430/}{391 U.S. at 439--41}).

\textbf{Holding vs. Dicta:} The \textit{holding} is that New Kent County's freedom-of-choice plan failed to satisfy the Board's affirmative remedial duty under \textit{Brown II} because after three years it left the dual system intact (zero white students at Watkins and 85\% of Black students still at Watkins) despite readily available alternatives like geographic zoning, and that the District Court must retain jurisdiction until state-imposed segregation is completely removed (\href{https://supreme.justia.com/cases/federal/us/391/430/}{391 U.S. at 439, 441--42}). By contrast, the Court's observations that ``there may well be instances in which [freedom of choice] can serve as an effective device'' (\href{https://supreme.justia.com/cases/federal/us/391/430/}{391 U.S. at 440--41}), its footnote recitation of the U.S. Commission on Civil Rights' findings on intimidation and poverty deterring Black transfers (which the Court expressly noted it ``neither adopt[s] nor refuse[s] to adopt,'' \href{https://supreme.justia.com/cases/federal/us/391/430/}{391 U.S. at 440 n.5}), and its discussion of specific alternative remedies such as consolidating the two schools by grade (grades 1--7 at one site and grades 8--12 at the other, \href{https://supreme.justia.com/cases/federal/us/391/430/}{391 U.S. at 442 n.6}) are \textit{dicta} offered as guidance for the remand rather than binding mandates for all districts.

\section{Analysis of Opinions}

\subsection{Majority Opinion (Justice Brennan)}

Justice Brennan began by framing New Kent County's system as the archetypal ``dual system, part `white' and part `Negro''' condemned by \textit{Brown I} and \textit{Brown II}, noting that racial identification extended across six operational dimensions: student body composition, faculty, staff, transportation, extracurricular activities, and facilities (\href{https://supreme.justia.com/cases/federal/us/391/430/}{391 U.S. at 435}). He rejected the School Board's core argument that adopting a facially race-neutral freedom-of-choice policy fully discharged its Fourteenth Amendment obligation and that anything more would improperly impose ``compulsory integration'' (\href{https://supreme.justia.com/cases/federal/us/391/430/}{391 U.S. at 437}). Opening the doors of formerly segregated schools in 1965 ``merely begins, not ends, our inquiry,'' because \textit{Brown II} imposed an affirmative duty on formerly \textit{de jure} segregated boards to eradicate all vestiges of state-imposed segregation ``root and branch'' (\href{https://supreme.justia.com/cases/federal/us/391/430/}{391 U.S. at 437--38}).

Timing and results were central to Brennan's reasoning. New Kent County had delayed taking even its first token step until eleven years after \textit{Brown I} and ten years after \textit{Brown II}, compounding the constitutional harm. Quoting \textit{Griffin v. County School Board of Prince Edward County}, 377 U.S. 218, 234 (1964), Brennan declared that ``[t]he time for mere `deliberate speed' has run out,'' shifting the standard from gradualism to immediate effectiveness: a board must present a plan that ``promises realistically to work, and promises realistically to work now'' (\href{https://supreme.justia.com/cases/federal/us/391/430/}{391 U.S. at 438--39}). District courts must independently evaluate proposed plans against feasible alternatives and retain jurisdiction until state-imposed segregation is completely eliminated (\href{https://supreme.justia.com/cases/federal/us/391/430/}{391 U.S. at 439}).

While declining to declare freedom-of-choice plans unconstitutional per se, Brennan adopted Judge Sobeloff's formulation from \textit{Bowman}: freedom of choice is ``not a sacred talisman,'' only a means to an end (\href{https://supreme.justia.com/cases/federal/us/391/430/}{391 U.S. at 440}). Where other reasonably available methods---such as geographic attendance zoning in a residentially integrated county, or splitting the two schools into one elementary school and one high school---promise speedier and more effective desegregation, freedom of choice is unacceptable (\href{https://supreme.justia.com/cases/federal/us/391/430/}{391 U.S. at 441--42 \& n.6}). Because New Kent's plan left 85\% of Black children at Watkins and attracted zero white students there, it improperly shifted the burden of desegregation from the School Board onto Black children and their families (\href{https://supreme.justia.com/cases/federal/us/391/430/}{391 U.S. at 441--42}).

\subsection{Concurrences}

No concurring opinions were filed.

\subsection{Dissents}

No dissenting opinions were filed.

\section{Examples: Future Applications}

\begin{enumerate}
\item \textbf{Same-Side Application (Free-Transfer Plan Preserving Dual Schools):} A formerly \textit{de jure} segregated city school district draws facially neutral geographic attendance zones that would integrate its schools, but adds a ``free-transfer'' provision allowing any student to transfer freely to any other school with available space. White students zoned into formerly Black schools overwhelmingly transfer back to formerly white schools, leaving the formerly Black schools nearly 100\% Black. \textit{Why Same-Side:} Like New Kent County's freedom-of-choice plan, the free-transfer option operates to perpetuate a dual system and shifts the burden of desegregation onto families when straightforward geographic zoning without racialized transfer escapes would work immediately (matching the companion decision in \href{https://caselaw.findlaw.com/court/us-supreme-court/391/450.html}{\textit{Monroe v. Board of Commissioners}}, 391 U.S. 450 (1968)).

\item \textbf{Same-Side Application (Residual Segregation Across Non-Student \textit{Green} Factors):} A rural county that previously operated statutory white and Black schools adopts geographic attendance zones that achieve rough racial balance in student enrollment, but continues assigning 95\% of Black teachers and support staff to the formerly Black campus, maintains inferior athletic facilities there, and runs racially distinct extracurricular clubs. \textit{Why Same-Side:} \textit{Green} requires eliminating racial identifiability ``root and branch'' across all six operational facets---not student assignment alone---so the district remains non-unitary and under continuing federal judicial supervision (\href{https://supreme.justia.com/cases/federal/us/391/430/}{391 U.S. at 435, 438}).

\item \textbf{Opposite-Side Application (Purely \textit{De Facto} Racial Imbalance Without Prior Intentional Segregation):} A northern suburban school district that never enacted or practiced intentional racial segregation assigns students strictly to neighborhood schools. Over two decades, private housing choices and economic shifts produce one elementary school that is 80\% Latino and Black and another that is 80\% white, with equal facilities, integrated faculty, and identical funding. \textit{Why Opposite-Side:} \textit{Green}'s affirmative duty to dismantle a dual system ``root and branch'' is triggered by prior state-imposed (\textit{de jure}) segregation; absent intentional segregative state action, racial imbalance alone does not violate the Equal Protection Clause (\href{https://columbialawreview.org/content/the-contested-role-of-time-in-equal-protection/}{Columbia Law Review}).

\item \textbf{Opposite-Side Application (Post-Unitary District Experiencing Demographic Resegregation):} A formerly \textit{de jure} segregated district complies in good faith with a federal desegregation decree for twenty years, satisfies all six \textit{Green} factors to the extent practicable, and obtains a federal court declaration of unitary status dissolving the injunction. Five years later, the board adopts a race-neutral neighborhood school policy that results in several predominantly single-race schools due to shifting residential demographics. \textit{Why Opposite-Side:} Once the vestiges of prior \textit{de jure} segregation have been eliminated to the extent practicable and unitary status is achieved, \textit{Green}'s remedial duty terminates and local control is restored unless plaintiffs prove new intentional discrimination (\href{https://scholarship.law.unc.edu/cgi/viewcontent.cgi?article=3435&context=nclr}{North Carolina Law Review}; \href{https://www.usccr.gov/files/historical/2006/06-006.pdf}{U.S. Commission on Civil Rights}).

\item \textbf{Fence-Sitter Application (Magnet-School Choice Plan With Partial Integration):} A formerly \textit{de jure} segregated county with moderate residential segregation replaces mandatory busing with a controlled-choice magnet program that places specialized STEM and arts academies at formerly Black schools, provides free countywide transportation, and pairs the plan with strict faculty integration. After three years, white enrollment at the formerly Black schools rises from 0\% to 30\% (in a 50/50 district), though mandatory cross-county reassignment could achieve 50/50 ratios immediately at the cost of substantial white flight. \textit{Why Fence-Sitter:} Under \textit{Green}, choice is not per se unconstitutional if it ``offers real promise'' and proves effective in practice (\href{https://supreme.justia.com/cases/federal/us/391/430/}{391 U.S. at 440--41}). A court must weigh whether 30\% integration with voluntary buy-in represents meaningful, realistic progress toward a unitary system or whether the availability of a faster mandatory zoning/pairing remedy renders the choice plan constitutionally inadequate.
\end{enumerate}

\section{Critique}

\textbf{Originalist and Textualist Critiques:} Conservative and originalist scholars have argued that \textit{Green} marked a pivotal---and textually ungrounded---shift from \textit{Brown I}'s colorblind prohibition against state-mandated racial classification to a mandate of race-conscious ``compulsory integration'' (\href{https://digitalcommons.law.seattleu.edu/cgi/viewcontent.cgi?article=1840&context=sulr}{Seattle University Law Review}). Critics contend that the Fourteenth Amendment's Equal Protection Clause forbids states from excluding students based on race, not from permitting genuinely voluntary parental choice; by measuring compliance through racial enrollment ratios, \textit{Green} invited federal district judges to engage in decades-long structural management of local schools, busing, and racial balancing that exceeded Article III remedial limits and strained the distinction between \textit{de jure} and \textit{de facto} segregation (\href{https://digitalcommons.law.seattleu.edu/cgi/viewcontent.cgi?article=1840&context=sulr}{Seattle University Law Review}; \href{https://www.usccr.gov/files/historical/2006/06-006.pdf}{U.S. Commission on Civil Rights}).

\textbf{Progressive and Doctrinal Perspectives:} Progressive scholars and civil rights advocates praise \textit{Green} as the decision that finally gave \textit{Brown} teeth, ending more than a decade of Southern ``Massive Resistance'' and tokenism by demanding actual results rather than paper neutrality (\href{https://tminstituteldf.org/renewing-the-promise-of-brown/}{LDF Thurgood Marshall Institute}; \href{https://socialchangenyu.com/harbinger/whatever-happened-to-the-green-factors-affirmative-action-through-the-lens-of-desegregation-law/}{NYU Review of Law \& Social Change}). At the same time, progressive commentators critique the doctrine's structural limits: because \textit{Green} tethered its affirmative remedial mandate strictly to remedying prior \textit{de jure} dual systems, later Supreme Court decisions confined \textit{Green} to districts with proven intentional segregation, insulated suburban districts from interdistrict remedies, and allowed federal courts to dissolve desegregation decrees once temporary ``unitary status'' was declared---contributing to widespread modern resegregation (\href{https://columbialawreview.org/content/the-contested-role-of-time-in-equal-protection/}{Columbia Law Review}; \href{https://scholarship.law.unc.edu/cgi/viewcontent.cgi?article=3435&context=nclr}{North Carolina Law Review}).

\textbf{Underlying Competing Values:} \textit{Green} sits at the core tension of Equal Protection jurisprudence: \textit{anticlassification} (formal race-neutrality, individual choice, and local democratic control of public schools) versus \textit{antisubordination} (substantive dismantling of entrenched racial caste systems and their lingering institutional effects through affirmative, results-oriented judicial remedies).

\section{Key Quotations}

\begin{quote}
``School boards such as the respondent then operating state-compelled dual systems were nevertheless clearly charged with the affirmative duty to take whatever steps might be necessary to convert to a unitary system in which racial discrimination would be eliminated root and branch.''
\end{quote}
\noindent --- Justice William J. Brennan, Jr., writing for a unanimous Court, \href{https://supreme.justia.com/cases/federal/us/391/430/}{391 U.S. at 437--38} (\href{https://www.courtlistener.com/opinion/107705/green-v-county-school-board-of-new-kent-county/}{CourtListener}).

\begin{quote}
``The burden on a school board today is to come forward with a plan that promises realistically to work, and promises realistically to work now.''
\end{quote}
\noindent --- Justice William J. Brennan, Jr., \href{https://supreme.justia.com/cases/federal/us/391/430/}{391 U.S. at 439} (\href{https://www.courtlistener.com/opinion/107705/green-v-county-school-board-of-new-kent-county/}{CourtListener}).

\begin{quote}
````Freedom of choice'' is not a sacred talisman; it is only a means to a constitutionally required end -- the abolition of the system of segregation and its effects. If the means prove effective, it is acceptable, but if it fails to undo segregation, other means must be used to achieve this end.''
\end{quote}
\noindent --- Justice William J. Brennan, Jr., quoting Judge Sobeloff in \textit{Bowman v. County School Board}, 382 F.2d at 333, at \href{https://supreme.justia.com/cases/federal/us/391/430/}{391 U.S. at 440} (\href{https://www.courtlistener.com/opinion/107705/green-v-county-school-board-of-new-kent-county/}{CourtListener}).

\begin{quote}
``Rather than further the dismantling of the dual system, the plan has operated simply to burden children and their parents with a responsibility which \textit{Brown II} placed squarely on the School Board. The Board must be required to formulate a new plan and, in light of other courses which appear open to the Board, such as zoning, fashion steps which promise realistically to convert promptly to a system without a `white' school and a `Negro' school, but just schools.''
\end{quote}
\noindent --- Justice William J. Brennan, Jr., \href{https://supreme.justia.com/cases/federal/us/391/430/}{391 U.S. at 441--42} (\href{https://www.courtlistener.com/opinion/107705/green-v-county-school-board-of-new-kent-county/}{CourtListener}).

\end{document}
